\documentclass[12pt,a4paper]{article}

% ─── Paquetes Básicos y Estructura ──────────────────────────────────────────
\usepackage[utf8]{inputenc}
\usepackage[T1]{fontenc}
\usepackage[spanish,es-tabla]{babel}

\usepackage[top=3.5cm, bottom=2.5cm, left=3cm, right=2.5cm, headheight=2.5cm, headsep=0.4cm]{geometry} 

\usepackage{amsmath, amssymb}
\usepackage{graphicx}
\usepackage{array}
\usepackage{tabularx}
\usepackage{booktabs}
\usepackage{enumitem}
\usepackage{fancyhdr}
\usepackage{titlesec}
\usepackage{xcolor}
\usepackage{multirow}
\usepackage{parskip}
\usepackage{tcolorbox}
\usepackage{wasysym}      
\usepackage{pifont}       

% Elimina los recuadros de color de las URLs de forma definitiva
\usepackage[hidelinks]{hyperref}

% ─── Configuración de Colores Estilo UTA ────────────────────────────────────
\definecolor{utablue}{RGB}{0,51,102}
\definecolor{resultbox}{RGB}{240,248,255}

% ─── Configuración de Párrafos del Documento ────────────────────────────────
\setlength{\parindent}{0pt}
\setlength{\parskip}{5pt}

% ─── Encabezado Oficial ──────────────────────────────────────────────────────
\pagestyle{fancy}
\fancyhf{} 

\fancyhead[C]{%
	\begin{minipage}[c]{0.13\textwidth}
		\centering
	\end{minipage}%
	\begin{minipage}[c]{0.74\textwidth}
		\centering
		\fontsize{11}{13}\selectfont\bfseries UNIVERSIDAD TÉCNICA DE AMBATO\\[3pt]
		\fontsize{9}{11}\selectfont\bfseries FACULTAD DE INGENIERÍA EN SISTEMAS, ELECTRÓNICA E INDUSTRIAL\\[2pt]
		\fontsize{9}{11}\selectfont\bfseries CARRERA DE SOFTWARE\\[2pt]
		\fontsize{9}{11}\selectfont\bfseries CICLO ACADÉMICO: ENERO 2026 -- JULIO 2026
	\end{minipage}%
	\begin{minipage}[c]{0.13\textwidth}
		\centering
	\end{minipage}\\[-2pt]
}

% Desactivación total de la numeración de páginas
\fancyfoot{} 
\renewcommand{\headrulewidth}{1.2pt}
\renewcommand{\footrulewidth}{0pt}

\makeatletter
\let\ps@plain\ps@empty
\makeatother
\pagestyle{fancy}

% ─── Formato de Secciones ───────────────────────────────────────────────────
\titleformat{\section}[block]
{\normalfont\bfseries\large}
{\Roman{section}.}{0.5em}{}

\titleformat{\subsection}[block]
{\normalfont\bfseries\normalsize}
{\thesubsection}{0.5em}{}

% ─── Cajas de Resultados Matemáticos ────────────────────────────────────────
\tcbuselibrary{skins}
\newtcolorbox{resultbox}{
	colback=resultbox,
	colframe=utablue,
	boxrule=0.8pt,
	arc=3pt,
	left=6pt, right=6pt, top=4pt, bottom=4pt
}

% ─── Interfaces de Casillas de Selección (Checkboxes) ───────────────────────
\newcommand{\cbox}{\makebox[1.2em]{\(\square\)}}
\newcommand{\cboxX}{\makebox[1.2em]{\(\boxtimes\)}}

% ════════════════════════════════════════════════════════════════════════════
\begin{document}
	
	\thispagestyle{fancy}
	
	\begin{center}
		\vspace*{0.1cm}
		{\large\bfseries\underline{INFORME CASOS DE ESTUDIO}}
	\end{center}
	
	\vspace{0.3cm}
	
	\section{Datos de los Estudiantes}
	
	\begin{tabular}{@{}p{6.0cm} >{\raggedright\arraybackslash}p{9.0cm}@{}}
		\textbf{Tema:}                              & Casos de estudio - Modelos Tradicionales \\[6pt]
		\textbf{Unidad de Organización Curricular:} & Profesional\\[6pt]
		\textbf{Nivel y Paralelo:}                  & Cuarto -- ``A'' \\[6pt]
		\textbf{Alumnos participantes:}             & Amaguaña Villacrés Shirley Dayana \\[6pt]
		\textbf  									& Caraguay Pucha Edwin Patricio\\[6pt]
		\textbf{Asignatura:}                        & Metodologías Ágiles \\[6pt]
		\textbf{Docente:}                           & Mgs. Hernan Naranjo \\
	\end{tabular}
	
	\vspace{0.4cm}
	
	% ────────────────────────────────────────────────────────────────────────────
	\section{Informe casos de estudio}
	\subsection{Análisis Modelo Incremental}
	
	
	% ────────────────────────────────────────────────────────────────────────────
	
	
	\vspace{0.5cm}
	
	\subsection{Conclusiones}
	
	\begin{enumerate}[leftmargin=1.5cm, label=\textbf{\arabic*.}, itemsep=6pt]
		\item La distribución uniforme continua es el modelo adecuado cuando todos los valores dentro de un intervalo $[a, b]$ son igualmente probables, como ocurre con la posición aleatoria de las manecillas de un reloj detenido.
		
		\item El cálculo de probabilidades en la distribución uniforme continua se simplifica a una razón entre longitudes de intervalo, lo que facilita su aplicación sin necesidad de integración compleja cuando los límites son conocidos.
		
		\item El resultado obtenido, $P(0 \leq X \leq 35) \approx 58{,}33\%$, indica que existe una probabilidad mayor al 50\% de que el reloj se haya detenido en los primeros 35 minutos, lo cual es consistente con que dicho intervalo representa más de la mitad del rango total de 60 minutos.
	\end{enumerate}
	
	\vspace{0.3cm}
	
	\subsection{Recomendaciones}
	
	\begin{enumerate}[leftmargin=1.5cm, label=\textbf{\arabic*.}, itemsep=6pt]
		\item Antes de aplicar la distribución uniforme continua, verificar que el fenómeno estudiado cumpla la condición de equiprobabilidad en todo el intervalo, ya que de no cumplirse se debe utilizar otro modelo de distribución continua.
		
		\item El orden recomendado de cálculo es primero resolver a mano, luego comprobar con tablas y finalmente con software estadístico como Excel o Jamovi, nunca a la inversa, para fortalecer la comprensión conceptual del modelo.
		
		\item Expresar siempre el resultado final en las tres formas: fracción exacta, decimal aproximado y porcentaje, dado que cada representación aporta una perspectiva diferente sobre la magnitud de la probabilidad calculada.
	\end{enumerate}
	
	\vspace{0.3cm}
	
	\subsection{Referencias bibliográficas}
	
	\begin{enumerate}[label={[\arabic*]}, leftmargin=1.5cm, itemsep=6pt]
		\item S. Monroy Saldivar, \textit{Estadística y probabilidad}, 1ra ed. México: Grupo Editorial Éxodo, 2024. [En línea]. Disponible en:
		\url{https://0x102qb9l-y-https-elibro-net.itmsp.museknowledge.com/es/lc/uta/titulos/275763}
		\item L. García Barco y L. M. Montoya Conde, \textit{Estadística descriptiva y fundamentos de probabilidad con R, RStudio y Geogebra}, 1ra ed. Colombia: Universidad Colegio Mayor de Cundinamarca, 2025. [En línea]. Disponible en:
		\url{https://0x102qb9l-y-https-elibro-net.itmsp.museknowledge.com/es/ereader/uta/291822}
		\item A. M. Grisales Aguirre, \textit{Estadística descriptiva y probabilidad con aplicaciones en EXCEL y SPSS}, 1ra ed. Colombia: Ecoe Ediciones, 2019. [En línea]. Disponible en:
		\url{https://0x102qb9l-y-https-elibro-net.itmsp.museknowledge.com/es/ereader/uta/125755}
		\item S. Monroy Saldivar, \textit{Estadística y probabilidad}, 1 ed. México: Grupo Editorial Éxodo, 2024. [En línea]. Disponible en:
		\url{https://elibro.net/es/ereader/uta/275763?page=161}
	\end{enumerate}
	
	\vspace{0.3cm}
	
	\subsection{Anexos}
	
	No aplica.
	
\end{document}